\documentclass[9pt]{extarticle}

\usepackage[margin=1in]{geometry}
\usepackage{amsmath,amssymb,mathtools,bm}
\usepackage[hidelinks]{hyperref}
\usepackage{microtype}

\newcommand{\R}{\mathbb{R}}
\newcommand{\diag}{\operatorname{Diag}}
\newcommand{\prox}{\operatorname{prox}}
\newcommand{\ones}{\bm{1}}
\newcommand{\had}{\mathbin{\odot}}

\title{Log-Barrier ADMM for Convex Quadratic Programs\\
\large A short note on smooth slack updates and implicit complementarity}
\author{}
\date{}

\begin{document}
\maketitle
\vspace{-2.5em}

\begin{abstract}
This note replaces the nonnegative-orthant indicator in a standard slack-variable
quadratic program by a logarithmic barrier and applies two-block ADMM.  The slack
subproblem has a closed-form proximal update.  Together with the subsequent dual
update, it produces exactly the paired retraction used by Arrizabalaga, and
Manchester~\cite{arrizabalaga2026implicit,arrizabalaga2026}, up to the natural ADMM scaling of the dual
variable and barrier parameter.  We also distinguish this elementary QP splitting
from the broader ADMM-based interior-point (ABIP) literature.
\end{abstract}

\section{Convex QP preliminaries}

Following~\cite{arrizabalaga2026,arrizabalaga2026implicit}, consider the convex quadratic program
\begin{equation}
\begin{aligned}
  \underset{x\in\R^n}{\operatorname{minimize}}\quad
      & \tfrac12 x^\top Qx+q^\top x \\
  \operatorname{subject\ to}\quad
      & Ax=b, \qquad Gx\le h,
\end{aligned}
\qquad
\Longleftrightarrow
\qquad
\begin{aligned}
  \underset{x\in\R^n,\,s\in\R^p}{\operatorname{minimize}}\quad
      & \tfrac12 x^\top Qx+q^\top x \\
  \operatorname{subject\ to}\quad
      & Ax=b,\qquad Gx+s=h,\qquad s\ge0,
\end{aligned}
\label{eq:qp}
\end{equation}
where $Q\in\mathbb{S}_+^n$, $q\in\R^n$, $A\in\R^{m\times n}$,
$b\in\R^m$, $G\in\R^{p\times n}$, and $h\in\R^p$.  The variables
$y\in\R^m$ and $z\in\R^p$ denote the multipliers of $Ax=b$ and
$Gx+s=h$, respectively.  The KKT residuals and cone conditions are
\begin{align}
 r_t&:=Qx+q+A^\top y+G^\top z=0, &
 r_c&:=z\had s=0, \nonumber\\
 r_e&:=Ax-b=0, &
 r_i&:=Gx+s-h=0, & s,z&\ge0.                 \label{eq:kkt}
\end{align}
Here $\had$ denotes componentwise multiplication.  Equivalently, the slack
formulation has objective
\begin{equation}
  f(x)+I_{\R_+^p}(s),\qquad
  f(x):=\tfrac12x^\top Qx+q^\top x,
\end{equation}
subject to the two affine equalities in~\eqref{eq:qp}.

For a central-path parameter $\kappa>0$, replace the indicator by the extended-valued
log barrier
\begin{equation}
  \phi_\kappa(s):=
  \begin{cases}
     -\kappa\sum_{i=1}^p\log s_i,&s>0,\\
     +\infty,&\text{otherwise}.
  \end{cases}
\end{equation}
The resulting barrier subproblem is
\begin{equation}
\begin{aligned}
  \underset{x,s}{\operatorname{minimize}}\quad
     & f(x)-\kappa\sum_{i=1}^p\log s_i\\
  \operatorname{subject\ to}\quad
     & Ax=b,\qquad Gx+s=h.                         \label{eq:barrier-qp}
\end{aligned}
\end{equation}
Its stationarity condition with respect to $s$ is
$z-\kappa s^{-1}=0$, where the inverse is componentwise.  Hence every solution of
\eqref{eq:barrier-qp} satisfies
\begin{equation}
   z\had s=\kappa\ones,\qquad s,z>0,               \label{eq:central}
\end{equation}
which is precisely the perturbed complementarity condition.  Sending
$\kappa\downarrow0$ traces the usual central path under standard regularity
conditions.

\section{ADMM derivation and the closed-form slack step}

Stack the affine constraints as
\[
 \begin{bmatrix}A\\G\end{bmatrix}x+
 \begin{bmatrix}0\\I\end{bmatrix}s=
 \begin{bmatrix}b\\h\end{bmatrix}.
\]
Thus~\eqref{eq:barrier-qp} is a two-block convex problem in $(x,s)$.  Using
unscaled multipliers $(y,z)$ and one penalty $\rho>0$, its augmented Lagrangian is
\begin{align}
 \mathcal{L}_{\rho,\kappa}(x,s,y,z)
 ={}& f(x)-\kappa\sum_i\log s_i
   +y^\top(Ax-b)+z^\top(Gx+s-h) \nonumber\\
 &+\frac{\rho}{2}\lVert Ax-b\rVert_2^2
  +\frac{\rho}{2}\lVert Gx+s-h\rVert_2^2.          \label{eq:augLag}
\end{align}
The corresponding ADMM iteration is
\begin{align}
 x^{k+1}&=H_\rho^{-1}\!\left(
   -q-A^\top y^k-G^\top z^k
   +\rho A^\top b+\rho G^\top(h-s^k)\right),       \label{eq:xstep}\\
 w^{k+1}&:=h-Gx^{k+1}-z^k/\rho,                    \label{eq:wdef}\\
 s^{k+1}&=\arg\min_{s>0}\left\{
   -\kappa\sum_i\log s_i
   +\frac{\rho}{2}\lVert s-w^{k+1}\rVert_2^2\right\}, \label{eq:sprox}\\
 y^{k+1}&=y^k+\rho(Ax^{k+1}-b),                    \label{eq:yupdate}\\
 z^{k+1}&=z^k+\rho(Gx^{k+1}+s^{k+1}-h).            \label{eq:zupdate}
\end{align}
where
\begin{equation}
 H_\rho:=Q+\rho A^\top A+\rho G^\top G.
\end{equation}
We assume $H_\rho$ is nonsingular; for example, it is positive definite if
$\ker Q\cap\ker A\cap\ker G=\{0\}$.  With fixed $\rho$, its factorization can be
reused across iterations.

The slack step is separable.  For one component, with
$\mu:=\kappa/\rho>0$, first-order optimality gives
\begin{equation}
 -\frac{\kappa}{s_i}+\rho(s_i-w_i)=0
 \quad\Longleftrightarrow\quad
 s_i^2-w_i s_i-\mu=0.
\end{equation}
Only the positive root lies in the barrier domain, so
\begin{equation}
 \boxed{
 s^{k+1}=b_\mu(w^{k+1}),\qquad
 b_\mu(t):=\frac{t+\sqrt{t^2+4\mu}}{2},\qquad
 \mu=\frac{\kappa}{\rho},}
 \label{eq:prox}
\end{equation}
with $b_\mu$ applied componentwise.  In proximal notation,
\begin{equation}
 b_\mu(w)=\prox_{\mu(-\log)}(w).
\end{equation}
This is the standard closed-form proximal operator of the negative logarithm;
see, for example, Parikh and Boyd~\cite{parikh2014}.

\section{Exact relation to implicit complementarity}

The map in~\eqref{eq:prox} is exactly the retraction selected in
\cite{arrizabalaga2026,arrizabalaga2026implicit}.  It satisfies
\begin{equation}
 b_\mu(t)b_\mu(-t)=\mu,\qquad
 b_\mu(t)-b_\mu(-t)=t,                              \label{eq:identities}
\end{equation}
and
\begin{equation}
 b_\mu'(t)=\frac12\left(1+\frac{t}{\sqrt{t^2+4\mu}}\right)\in(0,1),
 \qquad b_\mu'(t)+b_\mu'(-t)=1.                    \label{eq:derivative}
\end{equation}
The connection is not merely formal.  Let the scaled inequality multiplier be
$u^k:=z^k/\rho$.  Using~\eqref{eq:wdef} in~\eqref{eq:zupdate} yields
\begin{equation}
 u^{k+1}=s^{k+1}-w^{k+1}=b_\mu(-w^{k+1}).           \label{eq:paired}
\end{equation}
Consequently, immediately after the slack and dual updates,
\begin{equation}
 \boxed{
 u^{k+1}=b_\mu(v^{k+1}),\qquad
 s^{k+1}=b_\mu(-v^{k+1}),\qquad
 v^{k+1}:=-w^{k+1},}
 \label{eq:paper-map}
\end{equation}
and therefore
\begin{equation}
 u^{k+1}\had s^{k+1}=\mu\ones
 \quad\Longleftrightarrow\quad
 z^{k+1}\had s^{k+1}=\kappa\ones.                 \label{eq:exact-comp}
\end{equation}
Equations~\eqref{eq:paper-map}--\eqref{eq:exact-comp} are the implicit
complementarity representation of~\cite{arrizabalaga2026,arrizabalaga2026implicit}, expressed in ADMM's
scaled-dual coordinates.  In the paper's unscaled notation one may equivalently set
$v_{\rm paper}=z-s$ once~\eqref{eq:exact-comp} holds, giving
$z=b_\kappa(v_{\rm paper})$ and $s=b_\kappa(-v_{\rm paper})$.

This correspondence has two immediate implications.  First, positivity and perturbed
complementarity are enforced by construction in every completed $(s,z)$ update, not
only at convergence.  Second, the local derivative in~\eqref{eq:derivative} is bounded,
which makes the slack map smooth and compatible with automatic differentiation.  This
does \emph{not} by itself prove that gradients through arbitrarily many unrolled ADMM
iterations are well conditioned; conditioning also depends on the linear solve, the
iteration count, and the fixed-point Jacobian.  Implicit differentiation of the
converged fixed point is another option.

For numerical evaluation, the algebraically equivalent branch
\begin{equation}
 b_\mu(t)=\frac{2\mu}{\sqrt{t^2+4\mu}-t}\qquad(t<0)
\end{equation}
avoids cancellation in $t+\sqrt{t^2+4\mu}$, as emphasized for the same retraction
in~\cite{arrizabalaga2026,arrizabalaga2026implicit}.

\section{Relation to published work and scope}

The ingredients above are established, although their juxtaposition gives a useful
compact QP identity.
\begin{itemize}
\item \textbf{Proximal log barriers.}
The operator $\prox_{\mu(-\log)}=b_\mu$ is a standard example in the proximal
operator literature~\cite{parikh2014}.  Thus the quadratic formula in
\eqref{eq:prox} is not new by itself.

\item \textbf{ADMM-based interior-point methods.}
Lin, Ma, Ye, and Zhang introduced ABIP for large-scale linear programming
\cite{lin2021}.  ABIP uses ADMM to solve barrier-penalized subproblems in a
primal-dual path-following method built around a homogeneous self-dual embedding.
Deng et al. extend and improve this framework for linear and general conic
optimization~\cite{deng2025}.  These papers establish a substantially broader solver
architecture, including infeasibility handling, path-following logic, convergence
analysis, and conic proximal subproblems.

\item \textbf{What is specific here.}
The present derivation applies ordinary two-block ADMM directly to the slack form of
a convex QP and isolates the elementary update~\eqref{eq:prox}.  Its most direct link
to~\cite{arrizabalaga2026,arrizabalaga2026implicit} is the exact scaled-dual identity
\eqref{eq:paper-map}.  This note should therefore be viewed as a transparent
specialization and synthesis, not as a claim that ``ADMM plus a log barrier'' or the
closed-form proximal map is novel.
\end{itemize}

For a fixed $\kappa>0$, standard two-block ADMM theory applies to
\eqref{eq:barrier-qp} under the usual existence and regularity assumptions.  To solve
the original QP, one should place this inner iteration inside an outer schedule
$\kappa_j\downarrow0$, warm-starting each barrier subproblem.  Changing $\kappa$ at
every ADMM step without controlling inner accuracy is instead an inexact
path-following heuristic and requires a separate convergence argument.

\begin{thebibliography}{9}


\bibitem{arrizabalaga2026}
J.~Arrizabalaga, K.~Tracy, and Z.~Manchester.
\newblock A differentiable interior-point method in single precision.
\newblock \emph{arXiv preprint arXiv:2605.17913}, 2026.
\newblock \url{https://arxiv.org/abs/2605.17913}.

\bibitem{arrizabalaga2026implicit}
J.~Arrizabalaga and Z.~Manchester.
\newblock Implicit primal-dual interior-point methods for quadratic programming.
\newblock \emph{arXiv preprint arXiv:2604.00364}, 2026.

\bibitem{parikh2014}
N.~Parikh and S.~Boyd.
\newblock Proximal algorithms.
\newblock \emph{Foundations and Trends in Optimization}, 1(3):127--239, 2014.
\newblock \url{https://doi.org/10.1561/2400000003}.

\bibitem{lin2021}
T.~Lin, S.~Ma, Y.~Ye, and S.~Zhang.
\newblock An ADMM-based interior-point method for large-scale linear programming.
\newblock \emph{Optimization Methods and Software}, 36(2--3):389--424, 2021.
\newblock \url{https://doi.org/10.1080/10556788.2020.1821200}.

\bibitem{deng2025}
Q.~Deng, Q.~Feng, W.~Gao, D.~Ge, B.~Jiang, Y.~Jiang, J.~Liu, T.~Liu, C.~Xue,
Y.~Ye, and C.~Zhang.
\newblock An enhanced alternating direction method of multipliers-based interior
point method for linear and conic optimization.
\newblock \emph{INFORMS Journal on Computing}, 37(2):338--359, 2025.
\newblock \url{https://doi.org/10.1287/ijoc.2023.0017}.

\end{thebibliography}

\end{document}
